\chapter{Giấc Mộng Rớt Đại Học Vẫn Thành Tỷ Phú}
\markboth{Giấc Mộng Rớt Đại Học Vẫn Thành Tỷ Phú}{Giấc Mộng Rớt Đại Học Vẫn Thành Tỷ Phú}

\section{Lấy Ngoại Lệ Để Xóa Công Sức Số Đông}
\begin{truyen}{Lấy Ngoại Lệ Để Xóa Công Sức Số Đông}{Using Rare Exceptions to Erase the Majority's Effort}
\chuhoa{H}{ọc sinh:} ``Rớt đại học chưa chắc là hết đường. Thiếu gì người bỏ học vẫn thành công lớn.''
\textit{(Student: ``Failing university is not the end. Many dropouts still become hugely successful.'')}

\teacherpair{Câu đó đúng ở mức độ thống kê có tồn tại ngoại lệ. Nhưng đem ngoại lệ hiếm hoi ra để ru ngủ bản thân khỏi việc cố gắng là kiểu tự lừa rất sang mồm. Đa số người thành công nhờ lao động, học hỏi và tích lũy chứ không nhờ phép màu may mắn từ vài câu chuyện viral.}{``That is true only in the sense that exceptions exist. But using rare exceptions to lull yourself away from effort is a very polished form of self-deception. Most successful people grow through work, learning, and accumulation, not by magical luck drawn from a few viral stories.''}
\end{truyen}

\section{Mơ Lớn Không Sai, Mơ Tắt Nền Mới Nguy}
\begin{truyen}{Mơ Lớn Không Sai, Mơ Tắt Nền Mới Nguy}{Big Dreams Are Fine; Dreaming Without Foundations Is Dangerous}
\chuhoa{H}{ọc sinh:} ``Em đâu có chống học. Em chỉ nghĩ đâu cần học theo cách truyền thống.''
\textit{(Student: ``I am not against learning. I just think we do not always need the traditional route.'')}

\teacherpair{Không ai bắt em phải sống một bản sao. Nhưng bỏ một con đường chỉ đáng khi em thật sự đang xây con đường khác bằng kỷ luật, năng lực và sự nhất quán. Còn nếu bỏ nền tảng chỉ vì ngại khó, rồi gọi đó là tư duy khác biệt, thì đó không phải sáng tạo. Đó là né công việc nặng.}{``No one requires you to live as a copy. But leaving one path is worthwhile only if you are truly building another through discipline, ability, and consistency. If you abandon foundations merely because they are hard and call that unconventional thinking, that is not creativity. It is avoidance of heavy work.''}
\end{truyen}

\section{Mơ Giàu Thì Cũng Phải Tỉnh}
\begin{truyen}{Mơ Giàu Thì Cũng Phải Tỉnh}{If You Dream of Wealth, Stay Awake}
\chuhoa{H}{ọc sinh:} ``Vậy có được phép nghĩ lớn không thầy?''
\textit{(Student: ``Then am I allowed to think big, teacher?'')}

\teacherpair{Rất nên. Nhưng nghĩ lớn mà không chịu làm những việc nhỏ một cách tử tế thì cái lớn đó chỉ là bong bóng tự ái. Giấc mơ đáng tôn trọng luôn đi cùng kỷ luật buồn tẻ. Mơ giàu mà lười học, lười làm, lười sửa mình thì chỉ đang mơ cho đỡ chán thôi.}{``You absolutely should. But thinking big without doing small things properly turns ambition into an ego balloon. Dreams worth respecting always come with boring discipline. If you dream of wealth while refusing to learn, work, or correct yourself, you are only dreaming to entertain yourself.''}
\end{truyen}

\begin{baihoc}
Ngoại lệ là để truyền cảm hứng, không phải để hợp pháp hóa sự buông xuôi.
\textit{(Exceptions should inspire, not legalize surrender.)}
\end{baihoc}

\begin{ghinhoanh}
\textbf{Short English to remember:}

Dream big, but build small.

Rare stories are not safe plans.
\end{ghinhoanh}

\begin{tuhocnhanh}
1. Em đang mơ lớn thật, hay đang dùng mơ lớn để che nỗi sợ phải cố gắng? \textit{(Are you truly dreaming big, or using big dreams to cover your fear of effort?)}

2. Con đường khác biệt của em đã có nền thật chưa? \textit{(Does your unconventional path already have a real foundation?)}
\end{tuhocnhanh}

\ngancach
